% ECONOMICS_PROBLEM: {"id": "G11-625", "title": "Household stock-market nonparticipation and underdiversification", "jel_code": "G11", "source_id": "OP-0094"}
% !TeX program = xelatex
\documentclass[11pt,a4paper]{article}
\usepackage[margin=23mm]{geometry}
\usepackage{fontspec}
\setmainfont{Latin Modern Roman}
\usepackage{amsmath,amssymb,mathtools}
\usepackage{xcolor,enumitem,booktabs,longtable,array,xurl,hyperref}
\definecolor{muted}{HTML}{53636C}
\hypersetup{colorlinks=true,urlcolor=blue,linkcolor=blue}
\setlength{\parindent}{0pt}
\setlength{\parskip}{5pt}
\newcommand{\R}{\mathbb R}
\newcommand{\E}{\mathbb E}
\newcommand{\N}{\mathbb N}
\newcommand{\ind}{\mathbf 1}
\newcommand{\Var}{\operatorname{Var}}
\newcommand{\Cov}{\operatorname{Cov}}
\newcommand{\supp}{\operatorname{supp}}
\DeclareMathOperator*{\argmin}{arg\,min}
\DeclareMathOperator*{\argmax}{arg\,max}
\newcommand{\entrymeta}[3]{{\sffamily\small\color{muted}\textbf{\texttt{#1}}\quad|\quad #2\par #3\par}}
\newcommand{\Problemref}[1]{\texttt{#1}}
\newcommand{\JELref}[2]{\texttt{#2} (JEL \texttt{#1})}
\begin{document}
\section*{Household stock-market nonparticipation and underdiversification}
\textbf{Problem ID:} \texttt{G11-625}\quad \textbf{JEL:} \texttt{G11}\par
% BEGIN ECONOMICS STATEMENT
\label{problem:OP-0094}
\label{atlas:prob:F4}

\noindent\textbf{Atlas ID:} F4; \textbf{original number:} 94; \textbf{field:} Finance. \textbf{Classification:} Editorial research formulation (theoretical/empirical); not a certified unsolved theorem.

\subsection*{Definitions and assumptions}
For household $h$ over dates $0,\ldots,H$, let $a_{ht}$ be wealth, $y_{ht}$ nonfinancial income and $d_{ht}\in\{0,1\}$ stock-market participation. Participation incurs monetary cost $\kappa_h\ge0$. Portfolio weights $w_{ht}$ belong to a stated borrowing- and short-sale-feasible set, with zero equity exposure when $d_{ht}=0$. Consumption $c_{ht}>0$ and wealth satisfy
\[
 a_{h,t+1}=(a_{ht}+y_{ht}-c_{ht}-\kappa_hd_{ht})R_{h,t+1}(w_{ht}),
\]
subject to nonnegative investable resources or specified borrowing limits. Preferences, perceived risks, background income risk and survival probabilities form model $\theta$; households maximize finite expected discounted utility under their information.

\subsection*{Formal research question}
Identify the components of $\theta$ determining nonparticipation and concentration from wealth, income, holdings and belief data. For a randomized intervention $z$ altering information, trust or transaction costs, define $d_{hH}(z)$ and a prespecified concentration statistic $K_{hH}(z)$. Estimate participation effects and concentration effects on a fixed initial cohort; report concentration for nonparticipants under a stated convention or separately identify principal-stratum effects. Determine which intervention combinations distinguish fixed costs, pessimistic beliefs, background risk and preference heterogeneity, and quantify resulting household welfare rather than only increases in equity holdings.

\subsection*{Interpretation and scope}
A sum of squared portfolio weights measures concentration only relative to the defined holdings universe; employer shares, housing and entrepreneurial wealth can materially change the relevant exposure. Conditioning on participation after treatment introduces selection and is not a randomized comparison. The budget equation uses a gross portfolio return $R_{h,t+1}$ and is a transparent timing convention, not a claim that all households can trade one aggregate risky asset. Transaction costs may vary with trade size and participation history; replace the fixed cost if data require it. Information interventions can alter several primitives, so mediation effects need additional identification assumptions. General equilibrium welfare requires endogenous returns and subsidy financing.

\subsection*{Source audit and status}
[S1] addresses stockholding, [S3] surveys household finance, and [S2] explicitly disambiguates a relevant published Vissing-Jørgensen (2002) paper; another relevant 2002 working paper exists. These sources establish evidence and mechanisms in particular populations. They do not certify the broad household-portfolio question as an unresolved theorem or imply nonparticipation is always inefficient.

\subsection*{References}

\begin{enumerate}[label={[S\arabic*]},leftmargin=*]

\item Michael Haliassos and Carol C. Bertaut (1995). \emph{Why Do So Few Hold Stocks?}. Economic Journal 105(432), 1110–1129. \url{https://doi.org/10.2307/2235407}.
\textit{Locator and role:} Publisher or author abstract / bibliographic record; Household stock-market participation and proposed explanations. Provides background or a benchmark result; does not establish that the editorial question remains unresolved in 2026.

\item Annette Vissing-Jørgensen (2002). \emph{Limited Asset Market Participation and the Elasticity of Intertemporal Substitution}. Journal of Political Economy 110(4), 825–853. \url{https://doi.org/10.1086/340782}.
\textit{Locator and role:} Publisher or author abstract / bibliographic record; Participant-specific Euler-equation evidence. The author had more than one relevant 2002 paper. This published title is an explicit editorial disambiguation; NBER Working Paper 8884 is another relevant candidate.

\item John Y. Campbell (2006). \emph{Household Finance}. Journal of Finance 61(4), 1553–1604. \url{https://doi.org/10.1111/j.1540-6261.2006.00883.x}.
\textit{Locator and role:} Publisher or author abstract / bibliographic record; Survey of household portfolio choices and mistakes. Provides background or a benchmark result; does not establish that the editorial question remains unresolved in 2026.

\end{enumerate}

\subsection*{Common conventions: Source mathematical conventions}

Unless an entry states otherwise, agent and firm sets are finite. State and action spaces are measurable, strategies and outcome maps are measurable, probability laws are specified, and expectations exist for the targets under discussion. Infinite-horizon discount factors lie in $(0,1)$ and discounted objectives are finite. Norms, tolerances and complexity input sizes must be specified for computational comparisons. Constants or primitives that appear locally are inputs, not empirical facts asserted by this catalogue.

\textbf{Identification.} A statistical model $m\in\mathcal M$ implies an observable law $P_m$ and target $\theta(m)$. At law $P$, its identified set is
\[
\Theta(P;\mathcal M)=\{\theta(m):m\in\mathcal M,\ P_m=P\}.
\]
Point identification holds when this set is a singleton, or equivalently when $P_{m_1}=P_{m_2}$ implies $\theta(m_1)=\theta(m_2)$ for all compatible models. Sharp bounds mean bounds attained, or approached when the boundary is not attained, by compatible models. Writing this set is a definition, not a solution: the substantive task is to establish informative restrictions and derive or estimate the set. An unrestricted-confounding model may give only trivial bounds.

\textbf{Inference.} Identification, estimability and valid uncertainty quantification are separate questions. Fix $\alpha\in(0,1)$. A confidence procedure $C_n$ over a declared law class $\mathcal P$ has asymptotic uniform coverage at level $1-\alpha$ when
\[
\liminf_{n\to\infty}\inf_{P\in\mathcal P}\Pr_P\{\theta(P)\in C_n\}\geq1-\alpha.
\]
For set-identified targets the coverage event must be defined separately, for example coverage of the full identified set. If an entry uses identification-indexed models instead of uniquely indexed laws, the same coverage convention is applied over those models. Pointwise limits do not establish uniform validity, and asymptotic validity does not guarantee finite-sample precision. Sample sizes, dependence structures and weak-identification sequences are part of each application.

\textbf{Counterfactuals.} $Y_i(a)$ is a potential outcome under a specified intervention, including its timing, implementation and versions. It is not $Y_i$ conditional on voluntarily choosing $a$. A full intervention vector permits interference; shorthand scalar treatment notation does not silently impose no interference. Assignment, selection, attrition, anticipation and measurement must be specified. A claim about an instrument's local effect is not automatically a population policy effect.

\textbf{Economic equilibrium.} A counterfactual model includes best responses, constraints, feasibility, market clearing, state transitions and expectations consistent with the stated information. When several equilibria exist, either a declared selector is an input or the target is set-valued. Comparisons across policy regimes must use the stated selection convention. Reduced-form relationships are not presumed invariant after an equilibrium-changing intervention.

\textbf{Optimization and welfare.} Optimization is over the stated feasible set. If attainment is not established, $\sup$ or $\inf$ and $\varepsilon$-optimal choices replace an asserted maximizer or minimizer. For example, with value $V=\sup_{a\in\mathcal A}W(a)<\infty$, an $\varepsilon$-optimal set is $\{a\in\mathcal A:W(a)\geq V-\varepsilon\}$ for $\varepsilon>0$. Benefits, costs, risk and health or environmental outcomes can be aggregated only using declared units and conversion weights. Interpersonal utility weights, the treatment of fiscal transfers and foreign profits, discounting, and the behavioral welfare criterion are explicit inputs. A comparative-static derivative additionally requires existence and appropriate differentiability of the selected counterfactual.

\textbf{Sources.} Local labels [S1], [S2], and so forth restart in every question. Locators identify the relevant abstract, model, section or result at the level actually checked; a bibliographic or abstract-level verification is not represented as inspection of an entire proof. Working papers are distinguished from later publications and changing versions. Source summaries are paraphrased. All URL checks and editorial formulations refer to the audit date stated above.
% END ECONOMICS STATEMENT
\end{document}
